\section{Finite cut certificates and action regret}
\label{supp:computed126}

\subsection{Exact constraints and their rational tangents}

Use the notation $K=L-\mu$,
$h=q-\mu x^2/2$, $f_i=h(x_i)$, and $g_i=h'(x_i)$ from
Theorem~\ref{thm:fullinterpolation125}. Every item of information
claimed for the compatibility class must be encoded before computing a
cut. Additional original-model bounds are kept separately as described
below. In particular, an uncertain
supporting slope has one variable $p_j$ shared by every node. A
different slope chosen for each node would describe a different
information class.

Suppose first $K>0$. On an adjacent interval, let
$d=x_{i+1}-x_i$, $s=g_{i+1}-g_i$, and
$w=f_{i+1}-f_i-g_id$. In addition to the affine restrictions
$-s\le0$ and $s-Kd\le0$, the two nonlinear restrictions are
\[
 c_3=s^2-2Kw\le0,\qquad
 c_4=s^2-2K(sd-w)\le0.
\]
In the coordinate order $(f_i,f_{i+1},g_i,g_{i+1})$, their gradients are
\begin{align*}
 \nabla c_3&=(2K,-2K,-2s+2Kd,2s),\\
 \nabla c_4&=(-2K,2K,-2s,2s-2Kd).
\end{align*}
The Hessian of each is twice the outer product of the linear
coefficient vector of $s$ with itself. The remaining terms are affine.
Thus these are convex quadratics and their displayed tangents are
global affine minorants.

For a support at $a_j$ with level $W_j$, its node restriction is
\begin{align*}
 c_{ij}
 &=\ell_j(x_i;p_j)+\Phi_K(v;-x_i,A-x_i)-f_i\le0,\\
 v&=p_j-\mu a_j-g_i,\qquad
 t={\rm proj}_{[-x_i,A-x_i]}(v/K).
\end{align*}
Its only nonzero gradient coordinates are
\[
 \partial_{f_i}c_{ij}=-1,\qquad
 \partial_{g_i}c_{ij}=-t,\qquad
 \partial_{p_j}c_{ij}=x_i-a_j+t.
\]
The formula follows by differentiating
$\Phi_K(v;\ell,r)=vt-Kt^2/2$. The maximizer is unique when
$K>0$, including at the clipping thresholds, and the derivative with
respect to $v$ is $t$. Every evaluation is rational for rational
inputs. The function is convex because $\Phi_K$ is a pointwise
maximum of affine functions of $v$; subtracting $f_i$ and adding
$\ell_j$ preserves convexity.

When $K=0$, impose $s=w=0$ as affine equalities. The support
function is $\Phi_0(v;\ell,r)=\max\{v\ell,vr\}$.
Choose either maximizing endpoint $t\in\{\ell,r\}$, including
either one in a tie. The same three displayed coordinates are a valid
rational subgradient. This handles the constant-curvature case
without division by $K$ or an unstated limiting argument.

Observed value and derivative intervals, slope intervals, the cap
$|g_i+\mu x_i|\le M$, and the integrated value boxes are affine
restrictions. They can be retained as exact rows from the outset.
No fitted derivative, altered preference, or support from another
state may enter this set without its own declared transport theorem
and complete charge.

\subsection{Proof of Theorem~\ref{thm:cuts126}}

For every convex constraint, the subgradient inequality gives
\[
 c_{i_r}(z)\ge c_{i_r}(v_r)+s_r^\top(z-v_r)
                  =A_rz-b_r.
\]
If $z\in P$, its right side is nonpositive. Therefore, for
nonnegative $y$, $H(z)+y^\top(Az-b)\le H(z)$ on $P$.
The expression on the left is affine in $z$, and its exact minimum
on the finite product box is
\[
 h_0-y^\top b+\sum_k
       \min\{(h+A^\top y)_k\underline z_k,
              (h+A^\top y)_k\overline z_k\}=B.
\]
Thus $B\le\inf_{z\in P}H(z)$. The exact interpolation theorem
identifies the latter infimum as $-U_{\rm full}(a,b)$. This proves
the first assertion without any optimization assumption on the cut
points or weights.

If $z^{\rm f}\in P$, then
$B\le\min_PH\le H(z^{\rm f})$. Reversing signs gives
\eqref{eq:cutgap126}. In particular, the only accuracy assertion
licensed by the two witnesses is the actual checked gap. A favorable
floating-point objective, a small printed stationarity residual, or
an approximate primal feasibility flag does not replace either witness.

An upper bound remains valid after taking its minimum with other
valid upper bounds on the same number. The common node set represents
the same function class for each requested contrast; hence the same
valid cuts may be reused with a different linear $H$. Their generation,
storage, rational checks, and all numerical optimization calls remain
part of the service's work. If common acquisition is supplied to
different services, its allocation must be fixed and charged equally.
This proves all remaining assertions.\hfill$\square$

The assertion about the exact compatibility extremum applies to the
raw bound $-B$ and to intersections with bounds valid throughout that
same compatibility class. A separately licensed original-model bound
may describe a smaller class. In particular, the saved number
$B_{\rm old}\le\min q$ implies $q(a)\ge B_{\rm old}$ for the
actual original objective. The implemented cheap comparator uses this
fact even when the global restriction $q(x)\ge B_{\rm old}$ has
not been encoded in $P$. Its intersection with $-B$ certifies the
actual original action contrast, but need not dominate the supremum
over the larger $P$. Its economic validity does not by itself prove
the sharpness or numerical accuracy of the complete-support program.
The full-class regret sandwich uses full-class upper contrasts;
an original-policy regret upper bound may instead use any bounds
valid for that same actual $q$.

\paragraph{Numerical dual proposals.}
For a minimization linear program with rows $Az\le b$, a library may
report nonpositive objective derivatives with respect to the right-hand
sides. Under that convention the proposed nonnegative weights are
the negatives of those reported values. This is an implementation
convention to be verified, not a premise of the theorem. An independent
checker reconstructs $A,b$ from the original constraint identifiers
and rational cut points, rejects negative weights, and recomputes $B$.
Any rationalized nonnegative weights are valid, whether or not they
are close to a numerical optimum. Exact box minimization pays for all
stationarity error introduced by that rationalization. Retaining the
largest certified $B$ prevents a poorer later proposal from weakening
the certificate. Numerical linear-program infeasibility is never
treated as proof that the original economic record is impossible.

\subsection{The limiting calculation and its precise assumptions}

\begin{proposition}[Convergence of exact nested outer relaxations]
\label{prop:outerlimit126}
Let $\mathcal B$ be a compact convex box, $P\ne\varnothing$, and
let each convex $c_i$ be finite and continuous on a neighborhood of
$\mathcal B$, with bounded selected subgradients there. Start with
any finite valid cut set. At iteration $n$, minimize the affine
$H$ exactly on the current outer polytope $P_n\supseteq P$,
obtaining $z_n$. Add a tangent cut for every $c_i$ at $z_n$,
without discarding earlier cuts. Then
\[
 \min_{P_n}H\uparrow\min_PH,
\]
and every cluster point of $z_n$ belongs to $P$ and attains its
minimum. The conclusion also holds when cuts are taken at points
$v_n$ with $\|v_n-z_n\|\to0$, under the same bounded-subgradient
condition.
\end{proposition}

\begin{proof}
Each tangent is valid on $P$, so the nonempty compact outer
polytopes are nested. Their minima increase and are bounded above
by $\min_PH$; write their limit as $\ell$. Take any convergent
subsequence $z_{n_k}\to z_*$. For every fixed iteration $r$,
all sufficiently late points belong to the closed $P_r$, so
$z_*\in P_r$. In particular, the cut added at any $v_{n_k}$
is satisfied at $z_*$. Hence
\[
 c_i(v_{n_k})+s_{i,n_k}^\top(z_*-v_{n_k})\le0.
\]
Continuity, $v_{n_k}\to z_*$, and bounded subgradients imply
$c_i(z_*)\le0$ for every $i$. Thus $z_*\in P$.
Continuity of $H$ gives $H(z_*)=\ell$, which, combined with
$\ell\le\min_PH$, proves equality and optimality. The same
argument applies to every convergent subsequence.
\end{proof}

No interior point or constraint qualification is required for this
limiting statement. It gives no finite iteration count and no
convergence certificate for an arbitrary capped numerical solve.
Deleting cuts, adding cuts for only a selected subset of constraints,
or accepting a numerical subproblem solution without its required
control changes its hypotheses. The implemented economic guarantee
must therefore remain the finite rational bound of
Theorem~\ref{thm:cuts126}. When no exact primal witness is available,
the distance to the sharp information bound remains uncertified.

\subsection{Proof of Theorem~\ref{thm:gridregret126}}

For any fixed original objective in the information class,
\[
 q(a)-\min_xq(x)=\max_x\{q(a)-q(x)\}.
\]
Taking a supremum over the class and restricting the inner maximization
to the finite grid gives
\[
 \Gamma(a)\ge\max_j\sup_q\{q(a)-q(x_j)\}
                  =\max_j u_j.
\]
The maximum and supremum commute because the grid is finite. There
is no interchange of an infimum and a supremum in this argument.

For the upper bound, let $x_q^*$ minimize $q$. If it is a boundary
point, it is already a grid node. Otherwise $q'(x_q^*)=0$.
There is a grid node $x_j$ at distance at most $h/2$, and smoothness
gives
\[
 q(x_j)-q(x_q^*)\le\frac L2|x_j-x_q^*|^2
                                      \le\frac{Lh^2}{8}.
\]
It follows that
$q(a)-q(x_q^*)\le\max_j u_j+Lh^2/8$, uniformly in $q$.
Taking the supremum proves \eqref{eq:gridsandwich126}.

To verify sharpness, take the domain $[0,h]$, the grid $\{0,h\}$,
$a=0$, and the singleton class
$q(x)=L(x-h/2)^2/2$. All parameters can be positive rational
numbers and $0<\mu\le L$. Exact endpoint values and endpoint
derivatives with total derivative increment $Lh$ identify this
quadratic within the $L$-smooth class: for every intermediate $x$,
the upper Lipschitz bound from zero and the lower bound from $h$
coincide. Then $u_0=u_1=0$ and $\Gamma(a)=Lh^2/8$.

For the cell calculation, the function
$x\mapsto q(x)-Lx^2/2$ is concave on $[0,A]$. For
$x=(1-t)x_j+t x_{j+1}$ this gives the lower secant bound
\[
 q(x)\ge (1-t)q(x_j)+tq(x_{j+1})
                     -\frac{Ld_j^2}{2}t(1-t).
\]
After subtracting from $q(a)$ and applying the two verified
contrast bounds, the right side of \eqref{eq:cellregret126}
dominates $q(a)-q(x)$ on the entire cell. Maximizing over all cells
and then all compatible $q$ proves $\Gamma(a)\le G(a)$.
Because $a$ is a grid node, one adjacent cell endpoint evaluates
to zero; hence $G(a)\ge0$.

For a cell, write $k_j=Ld_j^2/2>0$. Its objective is
\[
 \widehat u_j+
  (\widehat u_{j+1}-\widehat u_j+k_j)t-k_jt^2.
\]
This concave quadratic is maximized at the clipped rational point
in \eqref{eq:cellmax126}. Finally its value is at most
$\max\{\widehat u_j,\widehat u_{j+1}\}+Ld_j^2/8$.
Under the uniform accuracy premise it follows that
\[
 G(a)\le\max_j u_j+\epsilon+Lh^2/8
       \le\Gamma(a)+\epsilon+Lh^2/8.
\]
This proves the accuracy statement. None of these arguments uses a
new value observation at a computational grid node. The exact
interpolation theorem applies to any enlarged finite node set,
so introducing the corresponding unknown values and derivatives
preserves the same compatible original functions.\hfill$\square$

\subsection{Robust action selection, original-policy gain and paid stopping}

For every compatible $q$, the function
$a\mapsto q(a)-\min q$ is $M$-Lipschitz. Taking suprema proves
the same Lipschitz bound for $\Gamma$. Also
$q(a)-\min q-\mu a^2/2$ is convex; its supremum is convex.
Thus $\Gamma$ is $\mu$-strongly convex. Compactness supplies a
unique minimizing action $a^{\rm r}$, which is the robust action
for this information class, not necessarily the actual optimum
$a_q^*$. Choose a grid node $a_j$ within $h/2$ of $a^{\rm r}$.
For a minimizer $\widehat a$ of the computed $G$ over the grid,
\[
 \Gamma(\widehat a)\le G(\widehat a)\le G(a_j)
 \le\Gamma(a_j)+\epsilon+Lh^2/8
 \le\Gamma(a^{\rm r})+Mh/2+\epsilon+Lh^2/8.
\]
This proves \eqref{eq:robustaction126}. The accuracy premise
requires a verified gap for every used pairwise calculation; a
solver's nominal tolerance is insufficient. Without that premise,
the selected action retains its finite regret certificate $G$,
but the displayed near-minimax guarantee is not asserted.

For the actual unchanged future controller,
$J(a)=q(a)+\eta(a)$ with $0\le\eta(a)\le\delta$,
\[
 0\le J(a)-\min q\le\Gamma(a)+\delta\le G(a)+\delta.
\]
If there is a separately certified original-value interval
$[W_a,U_a]$, the implemented policy interval is
$[W_a,U_a+\delta]$. The separate width gate is therefore
$U_a-W_a+\delta\le\omega$; it cannot be deduced from the
level-invariant regret calculation.

For two distinct realized actions,
$J(b)-J(a)\le U_{\rm full}(a,b)+\delta$.
Replacing $U_{\rm full}$ by any of its verified upper bounds
and requiring a strictly negative result proves actual gain for
every compatible original objective and allowed future error.
At identical actions, the unchanged implemented policy has exactly
zero contrast; it must not be assigned a positive gain by charging
two different future-error allowances. A finite sequence of
strictly accepted replacements telescopes its actual gains.

Finally, every feasible alternative original policy has cost at
least the original optimum. Starting from the qualified incumbent,
its possible economic saving is at most
$\min\{K_{\rm old},G(a)+\delta\}$. Under the stated compulsory
cost and no-reuse premise, its resource bill is at least
$\lambda c_*$. The stopping inequality therefore makes every
allowed continuation weakly more expensive in total objective.
Equality licenses weak dominance only. Reuse across future decisions,
an additional uncharged oracle, or another preference changes that
comparison and is excluded from the assertion.
